% ===== §27 =====
\section{The Production of Culture}
\label{sec:production}

\noindent
The third part turns from what intimate culture is and how it is lived to how it is produced and how it is captured, and it begins, as historical materialism requires, with production. The phenomenology treated the couple's culture largely from within, as it is generated and interpreted by the two; this part sets that internal life in the field of material forces that condition it from without, and asks the questions the first part's materialist line prepared: under what forces and relations of production is culture generated, who owns and transmits it, and how does it function in the reproduction of an order? The account proceeds through three layers---the forces of production, the relations of production, and ideology---and closes on the sharp contrast between two ways a subject can be produced, which is where the political economy meets the ethics of the series.

\subsection*{The forces of production: the materiality of the medium}

Culture is generated in a medium, and the medium is material. What forms a culture can take, how far they can travel, who can participate in making them, are set in the first instance by the forces of production---by the material means available for giving culture a body. The history of culture is inseparable from the history of its media: the passage from oral to written, from manuscript to print, from print to the electronic and the digital, is at each step a transformation of what culture can be, because each medium reconstitutes what can be recorded, how widely it can be shared, and who has access to the making. This is not a background fact but a condition internal to cultural production, and it bears on intimate culture as much as on the culture of a people. The forms a couple can generate and sustain are conditioned by the media through which they live and communicate: whether they share a household or are separated by distance, whether their communication passes through their own bodies and voices or through the platforms of others, whether the record of their shared life is kept in forms they control or in forms owned by another. The materiality of the medium is the first determinant of what an intimate culture can be, and it is, as the last sections of this part will show, the first point at which an external owner can insert itself into the couple's cultural production.

\subsection*{The relations of production: ownership and the position of the producer}

Beyond the forces of production stand the relations of production, and they raise the questions of ownership and of the producer's position that are decisive for a political economy of culture. Who has the leisure to generate culture; whose cultural production is recognized and whose is not; to whom the products belong; who holds the authority to interpret and to license---these are questions of the relations within which culture is produced, and their answers vary with the mode of production as sharply as the answers to the corresponding questions about material goods. A culture produced under aristocratic patronage, a culture produced in the guild, a culture produced under the culture industry, differ not only in content but in their whole relation of production: in who owns the product, who is positioned as its producer, and who is excluded from production altogether. For intimate culture the question of ownership seems, at first, not to arise---the couple generate their culture and it is theirs---but the appearance is deceptive, and the deception is the theme of this part. The question of who owns the medium through which the couple's culture is produced and recorded, who owns the platforms on which their shared life is increasingly lived, who is positioned to draw value from the culture they generate, turns out to be as live for the intimate dyad as for any other cultural producer, and the illusion that intimate culture stands outside the relations of production is precisely the illusion under cover of which it is captured.

\subsection*{Ideology: interpellation and the production of the subject}

The third layer is ideology, and it is the layer at which the production of culture becomes the production of subjects. Althusser gave this its sharpest form in the concept of \emph{interpellation}. Ideology, for Althusser, is not a set of ideas floating above the economic base but a material practice, embodied in institutions he called ideological state apparatuses---among which he named, decisively for this paper, the family---and its fundamental operation is to constitute individuals as subjects by \emph{hailing} them, calling to them in a way that, in the turning-to-answer, makes them the subject the hail addressed \citep{althusser1971}. The one hailed becomes, in recognizing themselves as the one called, a subject fitted to the order that called them; and this interpellation is, at bottom, the mechanism by which the relations of production reproduce themselves, producing the subjects they require. Culture, seen through this layer, is not innocent even in its most intimate forms, for the scripts of love, of the couple, of the shared life are among the hailings by which individuals are constituted as the subjects a given order needs---the subjects who will desire what the order supplies, mark their bonds in the forms the market sells, and reproduce, in their most private world, the relations of production that produced them.

\subsection*{Two ways to produce a subject: interpellation and generation}

Here the political economy meets the ethics of the series, at a contrast that must be drawn sharply, for it is the hinge on which the whole third part turns. There are two ways a subject can be produced in and through culture, and they are opposites. The first is interpellation in Althusser's sense: the production of a \emph{fitted} subject, hailed into a pre-existing position, constituted from without to occupy a role the order requires, formed as the subject that the relations of production need. This is the production of a subject as an effect, a congealing of the individual into a determined position, and it is, in the terms the series has developed, an operation of \emph{power-over}: the subject is produced by a power exercised upon it, fixed into the shape the order dictates. The second way is the one the whole of this paper has described: the generation of a subject within a relation, the coming-to-be of the relational subject through the culture it and its partner generate together. This is not the hailing of an individual into a fixed position but the generation of a subject that did not exist before, formed not from without by an order but from within by a relation, in the reciprocal movement by which the culture the couple generates comes back to constitute the couple. This is an operation of \emph{power-to}: an enabling generation, not a fixing subjection. The distinction is the ethical center of the political economy that follows. Interpellation produces the subject the order needs, from outside, by power-over; endogenous intimate culture generates the subject the relation becomes, from within, by power-to. And the capture of intimate culture by capital, which the next sections analyze, is precisely the operation by which the second is bent into the service of the first---by which the couple's own endogenous generation is colonized and made to produce, under the appearance of a freely generated intimate world, the fitted subjects the order requires. The Hegelian line names the danger exactly: production is objectification, the going-out of the subject into a made world, and objectification becomes \emph{alienation} when the made world is turned against its maker, when what the couple generates is owned by another and returns to them as a power over them rather than a world of their own. It is to that alienation, in its specific form as the capture of intimate culture by capital, that the argument now turns.
